\documentclass[11pt, a4paper]{article}
\usepackage[utf8]{inputenc}
\usepackage[spanish]{babel}
\usepackage{amsmath}
\usepackage{amsfonts}
\usepackage{amssymb}
\usepackage{graphicx}
\usepackage{geometry}
\usepackage{booktabs} % Para tablas profesionales (toprule, midrule, bottomrule)
\usepackage{fancyhdr}
\usepackage{xcolor}
\usepackage{float} % Para control de figuras (H)
\usepackage{array} % Para usar p{} en tablas
\usepackage{enumitem} % Para mejor control de listas

% Configuración de márgenes y estilo
\geometry{a4paper, total={170mm,257mm}, left=20mm, top=20mm}
\setlength{\headheight}{14pt} % Ajuste para eliminar warning de fancyhdr
\pagestyle{fancy}
\fancyhead{}
\fancyfoot{}
\fancyhead[L]{\textbf{\footnotesize Sistema de Recuperación de Calor CMPC}}
\fancyhead[R]{\textbf{\footnotesize Informe Final de Ingeniería de Detalle}} % Título mejorado
\fancyfoot[C]{\thepage}
\renewcommand{\headrulewidth}{0.4pt}
\renewcommand{\footrulewidth}{0.4pt}
\setlength{\tabcolsep}{5pt} % Espacio entre columnas de tablas

% Título y Metadata
\title{\vspace{-1.5cm}\textbf{\huge Propuesta de Inversión y Diseño:}}
\author{\textbf{\LARGE Sistema de Recuperación de Calor de Compresores (HRS)}\\
\normalsize Cliente: Papeles Cordillera S.A. (CMPC Puente Alto)}
\date{\normalsize Diciembre de 2025}


\begin{document}

\maketitle
\thispagestyle{fancy}

\newpage
\tableofcontents
\newpage

\section{Resumen Ejecutivo}

El presente informe detalla el diseño, dimensionamiento y análisis de rentabilidad para la implementación de un Sistema de Recuperación de Calor (HRS) en la estación de compresores de CMPC Puente Alto.

\subsection{Resultados Financieros Clave}

El proyecto es una inversión de alto retorno, demostrando una rentabilidad excepcional y listo para la ejecución inmediata.

\begin{table}[H]
    \centering
    \caption{Métricas Financieras y Operacionales del Proyecto HRS}
    \begin{tabular}{l r p{5cm}} % Métrica (l), Valor (r), Observación (p{5cm} para que ajuste el texto)
        \toprule
        \textbf{Métrica} & \textbf{Valor} & \textbf{Observación} \\
        \midrule
        Potencia Térmica de Diseño & $622.0 \text{ kW}$ & Condición Normal de Operación (80\%) \\
        Inversión Total (CAPEX) & $\mathbf{\$ 46,255,000 \text{ CLP}}$ & Costo Total Puesto en Planta (EPC) \\
        Ahorro Neto Anual & $\mathbf{\$ 273,259,210 \text{ CLP}/\text{año}}$ & Ahorro de Gas Evitado \\
        \textbf{Periodo de Recuperación (Payback)} & $\mathbf{0.17 \text{ años} (\sim 2 \text{ meses})}$ & \color{red}\textbf{PRIORIDAD MÁXIMA} \\
        Valor Actual Neto (VAN, 10 años) & $\mathbf{\$ 1,787 \text{ Millones CLP}}$ & Flujo Positivo Consolidado \\
        Tasa Interna de Retorno (TIR) & $\mathbf{590.8\%}$ & Rentabilidad Excepcional \\
        \bottomrule
    \end{tabular}
\end{table}

\subsection{Justificación de la Capacidad de Diseño}
El sistema se dimensionó para la \textbf{Condición Normal de Operación (Demanda 3)}, la cual ocurre el \textbf{80\%} del tiempo y corresponde a una potencia térmica disponible de $\mathbf{622 \text{ kW}}$.

\begin{figure}[H]
    \centering
    \includegraphics[width=0.9\textwidth]{scenario_analysis.png}
    \caption{Distribución de Escenarios de Operación y Potencia Térmica. La línea de diseño de $\mathbf{622 \text{ kW}}$ se alinea con el Escenario 3 (S3).}
\end{figure}

\section{Diseño de Componentes Críticos}

\subsection{Selección del Intercambiador de Calor (HX)}

La selección se basó en la eficiencia, el costo y el bajo ensuciamiento. Se determinó que el \textbf{Intercambiador de Placas} es la opción óptima, cumpliendo con la restricción de aislar el circuito del compresor del sistema de agua industrial.

\begin{table}[H]
    \centering
    \caption{Comparación de Tecnologías de Intercambio de Calor}
    \begin{tabular}{l c c c}
        \toprule
        \textbf{Tecnología} & \textbf{Área Requerida ($\mathbf{\text{m}^2}$)} & \textbf{Coef. U ($\mathbf{\text{W}/\text{m}^2\cdot\text{K}}$)} & \textbf{Costo Relativo} \\
        \midrule
        \textbf{Placas (Recomendado)} & $\mathbf{4.87}$ & $\mathbf{4000}$ & $\mathbf{1.0\times}$ \\
        Carcasa y Tubos & $19.46$ & $1000$ & $4.0\times$ \\
        \bottomrule
    \end{tabular}
\end{table}

\begin{figure}[H]
    \centering
    \includegraphics[width=0.9\textwidth]{heat_exchanger_design.png}
    \caption{Perfil de Temperaturas (Contraflujo) y Comparación Tecnológica del HX. Se valida el uso de placas por su alta eficiencia y bajo requerimiento de área.}
\end{figure}

\subsection{Sistema Hidráulico y Almacenamiento Térmico}

\subsubsection{Bombas de Circulación}
Se seleccionó un sistema de bombeo modular de baja potencia (\textbf{OPEX bajo}) y redundancia.
El caudal total de diseño es de $\mathbf{21.72 \text{ m}^3/\text{h}}$ (para $\mathbf{622 \text{ kW}}$ y $\Delta T=25 \text{ K}$), distribuido en 6 unidades de bombeo para modular la operación según el compresor activo.

\begin{figure}[H]
    \centering
    \includegraphics[width=0.9\textwidth]{pump_selection_analysis.png}
    \caption{Curvas de Sistema y Bombas Seleccionadas. La curva de la bomba se cruza con el punto de diseño del sistema (TDH), validando la selección y la eficiencia hidráulica.}
\end{figure}

\subsubsection{Tanque de Acumulación}
El tanque de $\mathbf{6.16 \text{ m}^3}$ provee la inercia térmica crítica ($\mathbf{15 \text{ minutos}}$ de almacenamiento) y está dimensionado para mantener la \textbf{estratificación térmica}.

\begin{figure}[H]
    \centering
    \includegraphics[width=0.9\textwidth]{storage_tank_design.png}
    \caption{Esquemático y Validación de la Estratificación del Tanque. El diseño de $H/D$ asegura el rendimiento térmico del almacenamiento.}
\end{figure}

\subsection{Ingeniería Básica: Fundamento Hidráulico y Componentes Auxiliares}

Esta sección consolida el fundamento técnico de la red hidráulica y los equipos auxiliares para garantizar la operación segura y la vida útil del sistema.

\subsubsection{Análisis de Caída de Presión y Tuberías}

El dimensionamiento de la tubería principal (circuito compresores-HX-Tanque) se realiza utilizando la fórmula de **Darcy-Weisbach** para minimizar la Altura Dinámica Total (TDH) y, por ende, el consumo eléctrico de las bombas.

\begin{itemize}[noitemsep, topsep=2pt]
    \item \textbf{Diámetro Seleccionado:} DN 50 (2 pulgadas) para el tramo principal, manteniendo la velocidad de flujo en el rango $\mathbf{1.5 \text{ m/s} - 2.0 \text{ m/s}}$.
    \item \textbf{Pérdida de Carga Total:} La pérdida de carga calculada se encuentra por debajo de $\mathbf{0.8 \text{ bar}}$, lo que resulta en un TDH manejable por las bombas seleccionadas.
\end{itemize}

\subsubsection{Componentes Auxiliares y Seguridad}

Los siguientes componentes auxiliares son esenciales para la integridad y el mantenimiento del sistema cerrado, y están incluidos en el presupuesto de CAPEX:

\begin{table}[H]
    \centering
    \caption{Componentes Auxiliares Críticos para la Operación}
    \begin{tabular}{p{3.5cm} p{6cm} c}
        \toprule
        \textbf{Componente} & \textbf{Función Crítica} & \textbf{Ubicación} \\
        \midrule
        Vaso de Expansión & Absorber la expansión térmica del agua ($\mathbf{T} = 85^{\circ}\text{C}$). & Circuito Cerrado \\
        Válvulas de Seguridad & Protección contra sobrepresión del circuito cerrado. & Tubería de Alta $\Delta P$ \\
        Filtros Y / Cestas & Prevenir el ingreso de partículas (óxido, suciedad) al HX y bombas. & Antes de Bombas/HX \\
        Purgadores de Aire & Eliminar aire del sistema, previniendo cavitación y corrosión. & Puntos Altos del Circuito \\
        Válvulas de 3 Vías & Control de temperatura y caudal para el HX (lado compresor). & Cabezal de Compresores \\
        \bottomrule
    \end{tabular}
\end{table}


\section{Análisis de Riesgos y Rentabilidad}

\subsection{Validación del Balance Energético}
El sistema cerrado de recuperación de calor presenta un error de balance de solo $\mathbf{0.83\%}$ ($\mathbf{5.15 \text{ kW}}$), confirmando la robustez termodinámica del diseño.

\subsection{Costo de Capital (CAPEX)}

\begin{table}[H]
    \centering
    \caption{Detalle de la Inversión Inicial (CAPEX)}
    % Uso de p{} para que las columnas de texto se ajusten y las cifras se alineen correctamente
    \begin{tabular}{p{3.5cm} p{6cm} r}
        \toprule
        \textbf{Componente} & \textbf{Descripción} & \textbf{Costo Estimado [CLP]} \\
        \midrule
        Intercambiador de Calor & Placas (Núcleo del sistema) & $19,000,000$ \\
        Tuberías y Red Hidráulica & Incluye aislamiento, fittings y auxiliares & $4,500,000$ \\
        Instrumentación & Sensores de $\Delta T$ y $\Delta P$, PLC & $3,500,000$ \\
        Otros Equipos & Bombas ($\mathbf{3.25\text{M}}$), Tanque ($\mathbf{2.8\text{M}}$) & $6,050,000$ \\
        \midrule
        \textbf{Costo de Ingeniería y Ejecución} & Ingeniería, Instalación, Contingencia (10\%) & $\mathbf{13,205,000}$ \\
        \midrule
        \textbf{TOTAL CAPEX} & & $\mathbf{\$ 46,255,000}$ \\
        \bottomrule
    \end{tabular}
\end{table}

\section{Conclusiones y Recomendaciones}

El proyecto es una oportunidad de \textbf{alta rentabilidad y bajo riesgo} que permite a CMPC Puente Alto ahorrar $\mathbf{616,223 \text{ m}^3}$ de gas natural anualmente.

\begin{itemize}
    \item La \textbf{TIR de $\mathbf{590.8\%}$} y el \textbf{Payback de 2 meses} posicionan a este proyecto como financieramente irrefutable.
    \item La validación técnica completa y el cumplimiento de los criterios de diseño confirman su \textbf{Preparación para Implementación (Ready-for-Implementation)}.
\end{itemize}

\subsection{Diagrama de Instrumentación y Flujo (P\&ID)}
El diagrama P\&ID es el entregable final clave para la fase de construcción y operación.

\begin{figure}[H]
    \centering
    \includegraphics[width=0.9\textwidth]{system_pid_diagram.png}
    \caption{Diagrama P\&ID del Sistema de Recuperación de Calor (HRS) CMPC. Muestra la interconexión de todos los componentes y el alcance del proyecto (Límite de Batería).}
\end{figure}

\end{document}